\section{A common-encoder obstruction with an interior dual witness}
\label{sec:v10-membership}

The following example puts the task-coupling and randomization issues
in one experiment. It is not a difference between distinct source
families or between incompatible loss normalizations.

The source reveals $X=\theta\in\{1,2,3\}$ without noise. Task
$d\in\{1,2,3\}$ asks whether $\theta=d$, with binary output and
zero-one loss. The encoder has two labels, is not told $d$, and is
common to the nine rows $(d,\theta)$. The task is supplied to the
decoder. All full-history task baselines are zero. A task-specific
two-label encoder would answer its task exactly.

\begin{theorem}[Three tasks, two labels]\label{thm:v10-membership}
For the common encoder, the best deterministic worst-row risk is one,
the best private-randomized risk is $1/2$, and the best shared-seed
risk is $2/9$. The uniform probability on all nine rows is a dual
optimizer for the shared problem. Every row receives positive dual
weight, and all nine risks of the displayed optimal shared scheme
are equal.
\end{theorem}
\begin{proof}
A deterministic two-label encoder merges at least two source symbols.
For the membership task of either symbol, the decoder makes an error
on at least one of them. Thus its worst-row risk is one. Moreover,
under the uniform distribution on the nine rows every deterministic
design makes at least two errors: each of the two membership tasks
for a merged pair has an error. The uniform risk is at least $2/9$.
Freezing all encoder and decoder randomization gives this same lower
bound for every shared mixture of complete designs.

For the matching scheme choose a singleton $k$ uniformly from the
three symbols using the public seed. Encode whether $X=k$. If $d=k$,
answer the task exactly. If $d\ne k$, answer zero on the singleton
label; on the pair label answer one with probability $2/3$. For a
positive row $d=\theta$, error occurs when $k\ne d$ and the decoder
answers zero, with probability $(2/3)(1/3)=2/9$. For a negative row
$d\ne\theta$, error occurs only when $k$ is the third symbol and the
decoder answers one, with probability $(1/3)(2/3)=2/9$. This establishes
optimality and the asserted strictly positive dual witness. The
additional decoder coin may be fresh private randomness; it may also
be frozen into the public mixture of deterministic designs.

For a private encoder write $p_x=\mathbb P(M=1\mid X=x)$. For each
task $d$, the probability of answering one is an affine function
$F_d(p_x)$, since it is obtained by applying two decoder probabilities
to the two encoder labels. One $p_d$ lies between the other two values,
including the possibility of equality. An affine function cannot be
greater than $1/2$ at this middle value and less than $1/2$ at both
end values. Hence some row of task $d$ has error at least $1/2$.
Ignoring the message and making a fair binary guess attains $1/2$.
A common retained coin shared with the encoder would invalidate this
affine argument, which is exactly why it is a different resource.
\end{proof}

For a single task, independent randomization cannot improve a scalar
expected-loss minimum, as in the strategic-measure results of
\cite{YukselSaldi2017}. The worst-row criterion here is a maximum of
linear risks, not one linear risk. Convexification can therefore
strictly improve it. There is no contradiction with that classical
irrelevant-randomization principle.
